% from ega2-1-fr.tex
\oldpage[II]{156}

\begin{proposition}[8.1.7]
\label{II.8.1.7-fr}
Soient $Y$ un préschéma, $\mathcal{I}$ un Idéal quasi-cohérent de
$\mathcal{O}_Y$, $X=\operatorname{Proj}(\bigoplus_{n\geq0}\mathcal{I}^n)$ le
$Y$-schéma obtenu en faisant éclater $\mathcal{I}$. On a alors, pour tout
$n>0$, un isomorphisme canonique
\begin{equation}
\label{II.8.1.7.1-fr}
  \mathcal{O}_X(n)\xrightarrow{\sim}\mathcal{I}^n\mathcal{O}_X
  =\mathcal{I}_{n,X}
\tag{8.1.7.1}
\end{equation}
(cf. (\hyperref[0.4.3.5-fr]{0, 4.3.5})), et par suite
$\mathcal{I}^n\mathcal{O}_X$ est un $\mathcal{O}_X$-Module inversible très
ample si $n>0$.
\end{proposition}

La dernière assertion est immédiate, puisque $\mathcal{O}_X(1)$ est inversible
(\hyperref[II.3.2.5-fr]{3.2.5}) et très ample pour $Y$ par définition
(\hyperref[II.4.4.3-fr]{4.4.3} et \hyperref[II.4.4.9-fr]{4.4.9}). D'autre
part, par définition, l'image de $v_n$ n'est autre que
$\mathcal{I}^n\mathcal{S}$, et (\hyperref[II.8.1.7.1-fr]{8.1.7.1}) résulte donc
de l'exactitude du foncteur $\widetilde{\mathcal{M}}$
(\hyperref[II.3.2.4-fr]{3.2.4}) et de la formule
(\hyperref[II.3.2.4.1-fr]{3.2.4.1}).

\begin{corollary}[8.1.8]
\label{II.8.1.8-fr}
Sous les hypothèses de (\hyperref[II.8.1.7-fr]{8.1.7}), si $f:X\to Y$ est le
morphisme structural et $Y'$ le sous-préschéma fermé de $Y$ défini par
$\mathcal{I}$, le sous-préschéma fermé $X'=f^{-1}(Y')$ de $X$ est défini par
$\mathcal{I}\mathcal{O}_X$ (isomorphe canoniquement à $\mathcal{O}_X(1)$),
d'où une suite exacte canonique
\begin{equation}
\label{II.8.1.8.1-fr}
  0\to\mathcal{O}_X(1)\to\mathcal{O}_X\to\mathcal{O}_{X'}\to0.
\tag{8.1.8.1}
\end{equation}
\end{corollary}

Cela résulte de (\hyperref[II.8.1.7.1-fr]{8.1.7.1}) et de
(\hyperref[I.4.4.5-fr]{I, 4.4.5}).

\begin{env}[8.1.9]
\label{II.8.1.9-fr}
Sous les hypothèses de (\hyperref[II.8.1.7-fr]{8.1.7}), on peut préciser la
structure des $\mathcal{I}_{n,X}$. Remarquons en effet que l'homomorphisme
\[
  \widetilde{u}_{-1}:\mathcal{O}_X\to\mathcal{O}_X(-1)
\]
correspond canoniquement à une section $s$ de $\mathcal{O}_X(-1)$ au-dessus de
$X$, que nous appellerons la \emph{section canonique} (relative à
$\mathcal{I}$) (\hyperref[0.5.1.1-fr]{0, 5.1.1}). D'autre part, dans le
diagramme (\hyperref[II.8.1.5.5-fr]{8.1.5.5}), les flèches horizontales sont
des isomorphismes (\hyperref[II.3.2.7-fr]{3.2.7})~; en remplaçant dans ce
diagramme $h$ par $k$ et $k$ par $-1$, on obtient
$\widetilde{u}_k=1_k\otimes\widetilde{u}_{-1}$ ($1_h$ désignant l'identité de
$\mathcal{O}_X(h)$), autrement dit l'homomorphisme $\widetilde{u}_k$ n'est
autre que la \emph{multiplication tensorielle par la section canonique $s$}
(pour tout $k\in\mathbf{Z}$). L'homomorphisme $\widetilde{u}$
(\hyperref[II.8.1.5.6-fr]{8.1.5.6}) s'interprète donc aussi de la même manière.

Par suite, pour tout $n\geq0$, l'homomorphisme
$\widetilde{v}_n:\mathcal{O}_X(n)\to\mathcal{O}_X$ n'est autre que la
multiplication tensorielle par $s^{\otimes n}$~; on en déduit~:
\end{env}

\begin{corollary}[8.1.10]
\label{II.8.1.10-fr}
Avec les notations de (\hyperref[II.8.1.8-fr]{8.1.8}), l'espace sous-jacent à
$X'$ est l'ensemble des $x\in X$ tels que $s(x)=0$, $s$ désignant la section
canonique de $\mathcal{O}_X(-1)$.
\end{corollary}

En effet, si $c_x$ est un générateur de la fibre $(\mathcal{O}_X(1))_x$ en un
point $x$, $s_x\otimes c_x$ s'identifie canoniquement à un générateur de la
fibre de $\mathcal{I}_{1,X}$ au point $x$, et est donc inversible si et
seulement si l'on a
$s_x\notin\mathfrak{m}_x(\mathcal{O}_X(-1))_x$, autrement dit si
$s(x)\neq0$.

\begin{proposition}[8.1.11]
\label{II.8.1.11-fr}
Soient $Y$ un préschéma intègre, $\mathcal{I}$ un Idéal fractionnaire
quasi-cohérent de $\mathcal{R}(Y)$, $X$ le $Y$-schéma obtenu en faisant éclater
$\mathcal{I}$. Alors $\mathcal{I}\mathcal{O}_X$ est un
$\mathcal{O}_X$-Module inversible très ample pour $Y$.
\end{proposition}

La question étant locale sur $Y$ (\hyperref[II.4.4.5-fr]{4.4.5}), on peut se
borner au cas où $Y=\operatorname{Spec}(A)$, où $A$ est un anneau intègre de
corps des fractions $K$ et $\mathcal{I}=\widetilde{\mathfrak{I}}$,
$\mathfrak{I}$ étant un idéal fractionnaire de $K$~; il existe par suite un
élément $a\neq0$ de $A$ tel que $a\mathfrak{I}\subset A$. Posons
$S=\bigoplus_{n\geq0}\mathfrak{I}^n$~; l'application $x\mapsto ax$ est un
$A$-isomorphisme de $\mathfrak{I}^{n+1}=(S(1))_n$ sur
$a\mathfrak{I}^{n+1}=a\mathfrak{I}S_n\subset\mathfrak{I}^n=S_n$,

